% ---------------------------------------------------------------------
%  Sample content: numbered boxes
% ---------------------------------------------------------------------
\section{Numbered boxes}\label{sec:boxes}

Numbered boxes hold the content you will refer back to. Each one takes an
optional title and an optional label; the number follows the section.

\begin{definition}[Title of the definition][def:sample]
A \emph{definition} states precisely what a term means. For example, a
function $f$ is \emph{increasing} if $x < y$ implies $f(x) \le f(y)$.
\end{definition}

\begin{theorem}[Title of the theorem][thm:sample]
A \emph{theorem} states a main result. If $f$ and $g$ are increasing, then
$f + g$ is increasing.
\end{theorem}

\begin{proof}
Let $x < y$. By \cref{def:sample}, $f(x) \le f(y)$ and $g(x) \le g(y)$.
Adding the two inequalities gives $(f + g)(x) \le (f + g)(y)$.
\end{proof}

\begin{lemma}
Lemmas, propositions and corollaries share the color and the counter of
theorems.
\end{lemma}

\begin{corollary}
By \cref{thm:sample}, any finite sum of increasing functions is increasing.
\end{corollary}

\begin{example}[][ex:sample]
Examples are shown in blue. Both $f(x) = x$ and $g(x) = x^3$ are increasing,
so $x + x^3$ is increasing as well.
\end{example}

\begin{remark}
Remarks are gray, for side comments that do not need emphasis.
\end{remark}

\begin{exercise}[][exer:sample]
Exercises are orange. Show that $h(x) = -x$ is not increasing.
\end{exercise}

\begin{solution}
Take $x = 0 < 1 = y$. Then $h(0) = 0 > -1 = h(1)$. Solutions can be hidden
for handouts with \code{\showsolutionsfalse}.
\end{solution}
